\documentclass[12pt]{report}
\usepackage[utf8]{inputenc}
\usepackage{graphicx}
\graphicspath{ {images/} }
\usepackage[a4paper, total={6in, 8in}]{geometry}

%% Language and font encodings
\usepackage[english]{babel}
\usepackage{graphicx}
\graphicspath{ {images/} }
\usepackage{setspace}
\usepackage{ragged2e}


%% Sets page size and margins
\usepackage[a4paper,top=3cm,bottom=2cm,left=3cm,right=3cm,marginparwidth=1.75cm]{geometry}

%% Useful packages
\usepackage[utf8x]{inputenc}
\usepackage{titlesec}
\include{preamble.tex}
\include{bibliography}
\usepackage{expex}


\setcounter{secnumdepth}{4}


\titleformat{\paragraph}
{\normalfont\normalsize\bfseries}{\theparagraph}{1em}{}
\titlespacing*{\paragraph}
{0pt}{3.25ex plus 1ex minus .2ex}{1.5ex plus .2ex}


\title {{Subject focus in French and Spanish}
{\begin{figure}
    \centering
    \includegraphics[width=0.2\textwidth ]{{university.png}}
\end{figure}
\author{Alessia Cassarà}}}

\date{}


\begin{document}
\doublespacing
\maketitle
\newpage


\renewcommand*\contentsname{Table of contents} 
\tableofcontents
\newpage

\chapter*{General introduction}\addcontentsline{toc}{chapter}{Introduction}
\newpage
\section*{Why investigate subject focus?} \label{why}\addcontentsline{toc}{section}{Why investigate subject focus?}
\section*{Terminology} \addcontentsline{toc}{section}{Terminology}
\section*{Goals and structure of the dissertation}\addcontentsline{toc}{section}{Goals and structure of the dissertation}
\part{Theoretical Background}
\chapter{Focus and focus realization}
\section{The notion of Focus}
\subsection{The formal approaches}
\subsection{The operational approach}
\subsection{Focus structure}
\subsection{Focus types}
\section{Focus realization}
\subsection{The argument asymmetry in focus realization}
\chapter{Focus subjects in Spanish and French}
\section{Introduction}
\section{Subject focus in Spanish}
\subsection{Postverbal and preverbal subjects: formal and functional approaches}
\subsection{Alternative strategies for Spanish focused subjects}
\section{Subject focus in French}
\subsection{Clefts as focus-marking strategies}
\subsection{Alternative strategies for French focused subjects}
\section{Comparative studies on Spanish and French focused subjects}
\subsection{French clefts and Spanish postverbal subjects in comparison}
\subsection{The factors influencing both structures}
\section{Filling the gap}
\part{Corpus study}
\chapter{The realization of focus on the subject: a comparative corpus study}
\section{Introduction}
\section{The \textit{sgs}corpus}
\section{Methodology}
\subsection{Criteria for the extraction}
\subsection{Criteria for the annotation}
\section{Focus-marking strategies: results}
\subsection{Corpus results for Spanish}
\subsection{Discussion}
\subsection{Corpus results for French}
\subsection{Discussion}
\section{Comparative analysis of French and Spanish}
\section{Conclusion chapter 4}
\chapter{French clefts and Spanish subject position: applying the same predictors}
\section{Introduction}
\section{Hypothesis}
\section{Methodology}
\subsection{The predictors}
\subsection{The argument structure of the verb: criteria for annotation}
\subsection{Activation of the subject constituent: criteria for annotation}
\section{Results}
\subsection{Spanish: results}
\subsection{French: results}
\section{Discussion}
\section{Comparative analysis}
\section{Open issues}
\section{Conclusion chapter 5}
\part{The pragmatic account}
\chapter{Motivating the choice of different focus-marking strategies}
\section{Introduction}
\subsection{The concept of Economy of the language}
\subsubsection{Grice’s maxims}
\subsubsection{Common Ground management and activation of given material}
\section{Hypothesis}
\section{Methodology}
\subsection{Material: full forms vs. reduced forms}
\subsection{Criteria for the annotation}
\section{Results}
\subsection{Results for Spanish}
\subsection{Results for French}
\section{Discussion}
\section{Comparative analysis}
\section{Conclusion chapter 6}
\chapter*{General conclusion} \addcontentsline{toc}{chapter}{General conclusion}


\end{document}